% Part: lambda-calculus
% Chapter: syntax
% Section: eta

\documentclass[../../../include/open-logic-section]{subfiles}

\begin{document}

\olfileid{lam}{syn}{eta}
\olsection{$\eta$-રૂપાંતરણ}

$\lambda$-પદો પર બીજો પણ એક સંબંધ છે. \olref[fv]{sec}માં
આપણે $\lambd[x][(fx)]$નું ઉદાહરણ જોયું હતું: તે આર્ગ્યુમેન્ટ
સ્વીકારીને તેના પર~$f$ લાગુ કરે છે. બીજા શબ્દોમાં, તે~$f$
જેવું જ વિધેય છે: $\lambd[x][(fx)]N$ અને~$fN$ બંનેનું
લઘુકરણ~$fN$માં થાય છે. આ વિચાર વ્યક્ત કરવા
$\eta$-લઘુકરણ (અને $\eta$-વિસ્તરણ) વાપરીએ છીએ.

\begin{defn}[$\eta$-સંકોચન, $\eredone$]
  \ollabel{defn:beredone}
  \emph{$\eta$-સંકોચન} ($\eredone$) પદો પરનો સૌથી નાનો
  રચના-સુસંગત સંબંધ છે જે નીચેની શરત સંતોષે છે:
  \[
    \lambd[x][M x] \eredone M \text{ જો } x \notin FV(M)
  \]
\end{defn}

\begin{defn}[$\beta\eta$-લઘુકરણ, $\bered$] \ollabel{defn:bered}
  \emph{$\beta\eta$-લઘુકરણ} ($\bered$) પદો પરનો
  $\bredone$ અને $\eredone$ને સમાવતો સૌથી નાનો
  સ્વવાચક અને પરંપરિત સંબંધ છે. એટલે સ્વવાચકતા અને
  પરંપરિતતાના નિયમો ઉપરાંત નીચેના બે નિયમો છે:
  \begin{enumerate}
  \item જો $M \bredone N$, તો $M \bered N$. \ollabel{defn:bered3}
  \item જો $M \eredone N$, તો $M \bered N$. \ollabel{defn:bered4}
  \end{enumerate}
    
\end{defn}

\begin{defn}
  સામ્ય સંબંધ~$\equal$માં $\eta$-રૂપાંતરણનો નિયમ ઉમેરીએ:
  \[
  \lambd[x][f x] \equal f
  \]
  આ રીતે વિસ્તરેલા સંબંધને $\equal[\eta]$થી દર્શાવીએ.
\end{defn}

$\eta$-સામ્ય મહત્વનું છે, કારણ કે તે લૅમ્બડા પદોની
વિસ્તારાત્મકતા સાથે સંબંધિત છે:

\begin{defn}[વિસ્તારાત્મકતા]
  સામ્ય સંબંધ~$\equal$માં (\ext) નિયમ ઉમેરીએ:
  \begin{center}
  જો $Mx \equal Nx$, તો $M \equal N$; શરત છે કે $x \notin FV(MN)$.
  \end{center}
  આ રીતે વિસ્તરેલા સંબંધને $\equal[\ext]$થી દર્શાવીએ.
\end{defn}

આ નિયમનો સાર એ છે કે પદોને વિધેયો માનીએ ત્યારે એક જ
મનસ્વી આર્ગ્યુમેન્ટ પર તેમનું વર્તન સમાન હોય તો તેમને
સમાન ગણવાં.

હવે સાબિત કરીએ કે $\eta$-નિયમથી બરાબર વિસ્તારાત્મકતા
મળે છે, તેનાથી વધુ કંઈ નહીં.

\begin{thm}
  $M \equal[\ext] N$ ત્યારે અને માત્ર ત્યારે જ જ્યારે
  $M \equal[\eta] N$.
\end{thm}

\begin{proof}
  પ્રદર્શિત $\eta$-નિયમમાં $f$ ચલ છે. પરંતુ
  $x \notin FV(M)$ હોય અને~$f$ને $M$ તથા~$x$થી
  તાજો લઈએ તો રચના-સુસંગતતાથી
  $\lambd[f][\lambd[x][f x]] \equal[\eta]
  \lambd[f][f]$. બંને બાજુ~$M$ લાગુ કરી
  $\beta$-સામ્ય વાપરતાં $\lambd[x][Mx]
  \equal[\eta] M$ મળે છે. આથી આગળ સામાન્ય પદો માટે
  $\eta$-નિયમ વાપરી શકીએ.
  \sourcecorrection{OLLAM-030}{સ્થિર અંગ્રેજી મૂળના
  પ્રદર્શિત $\eta$-નિયમમાં $f$ ચલ છે, પરંતુ સાબિતી
  ગમે તે પદ~$M$ માટે એ નિયમ વાપરે છે. અહીં તાજો
  $f$, રચના-સુસંગતતા અને $\beta$-સામ્યથી સામાન્ય
  નિયમ મેળવવાનું ગાયબ પગલું ઉમેર્યું છે.}
  પહેલાં બતાવીએ કે $\equal[\eta]$ વિસ્તારાત્મકતાના
  નિયમ હેઠળ બંધ છે. તેથી~\ext{} નિયમ
  $\equal[\eta]$માં નવી જોડી ઉમેરતો નથી: નિયમ-ઉપજ
  પર આગમનથી $\equal[\ext]$નો દરેક સંબંધ
  $\equal[\eta]$માં આવે છે. તેથી
  $M \equal[\ext] N$ હોય તો $M \equal[\eta] N$.

  \ext{} નિયમ હેઠળની બંધતા માટે, એ નિયમ લાગુ કરતાં
  $Mx \equal[\ext] Nx$ એ પૂર્વધારણા હોય છે અને
  $x \notin FV(MN)$ હોય છે. આગમનની ધારણાથી
  $Mx \equal[\eta] Nx$. રચના-સુસંગતતાથી
  $\lambd[x][Mx] \equal[\eta] \lambd[x][Nx]$.
  $x$ બંને પદોમાં મુક્ત ન હોવાથી, બંને બાજુ
  $\eta$-નિયમ લગાડી પરંપરિતતાથી
  $M \equal[\eta] N$ મળે છે.
  \sourcecorrection{OLLAM-031}{સ્થિર અંગ્રેજી મૂળ
  \ext{}-નિયમની પૂર્વધારણાને તરત
  $Mx \equal[\eta] Nx$ કહે છે, જ્યારે તે શરૂઆતમાં
  $Mx \equal[\ext] Nx$ છે; પહેલાંના નિષ્કર્ષ પર
  આગમનની ધારણા લગાવવી પડે. અહીં એ પગલું અને
  $x \notin FV(MN)$ની શરત સ્પષ્ટ કર્યાં છે.
  મૂળના બીજા દિશાના તર્કમાં $\equal[ext]$ લખાયું
  છે; વ્યાખ્યાયિત સંકેત $\equal[\ext]$ છે.}

  ઊલટું, $\eta$-નિયમ $\equal[\ext]$માં સમાયેલો છે.
  જો $x \notin FV(M)$ હોય તો $\beta$-નિયમથી
  $(\lambd[x][Mx])x \equal[\ext] Mx$. વળી,
  $x \notin FV((\lambd[x][Mx])M)$; તેથી
  \ext{} નિયમથી $\lambd[x][Mx] \equal[\ext] M$.
  આ રીતે
  $\equal[\eta]$ના નિયમો પણ $\equal[\ext]$માં
  સમાયેલા છે; એટલે બંને સંબંધો સમાન છે.
\end{proof}

\end{document}
